% ============================================================
%  ch02 —— §1 结构群（原文 p.7–p.11）
% ============================================================
\routenote{p.7--p.11}

\section{结构群}

我们将在特征为 $p$ 的固定有限域 $k$ 上工作。固定底曲线 $\PP^1$ 的坐标 $t$，使得 $\Aff^1=\operatorname{Spec}(k[t])\subset\PP^1$。记 $s:=t^{-1}$ 为 $\infty\in\PP^1$ 附近的坐标。

对任意闭点 $x\in\PP^1$，记 $\mathcal O_x$ 为 $x$ 处的完全局部环，$K_x$ 为 $\mathcal O_x$ 的分式域。

由于我们关心的是（野）分歧局部系统的几何 Langlands 对应，我们需要考虑带各种水平结构的主丛。把这些主丛看作 $\PP^1$ 上群概形 $\mathbb{G}$ 的 torsor 会很方便。此外，为了陈述猜想中的函子性，允许拟分裂群概形是有益的。我们将分几步引入这些群概形。

\subsection{$\PP^1\setminus\{0,\infty\}$ 上的拟分裂群概形}

我们假设 $\mathbb{G}|_{\PP^1\setminus\{0,\infty\}}$ 是拟分裂约化群。此外假设存在有限扩张 $k'/k$ 与整数 $N$，满足 $(N,p)=1$，使得 $\mathbb{G}$ 在驯扩张 $[N]:\GG_{m,k'}\to\GG_{m,k'}$（由 $t\mapsto t^N$ 定义）上分裂。我们可以且将会假设 $k'$ 包含所有 $N$ 次单位根。记 $\mu_N$ 为 $N$ 次单位根之群。

固定子群 $S\subset T\subset B\subset G$，其中 $S$ 是极大分裂环面，$T$ 是极大环面，$B$ 是 Borel 子群。我们还固定 $G$ 的一个拟钉住（quasi-pinning）。

为描述我们的群概形，为简单起见假设 $k=k'$。一般情形下，我们给出的构造在 $\operatorname{Gal}(k'/k)$ 下不变，于是 Galois 下降即给出一般情形。

由 [13]（Exp.~XXIV, Thm~3.11），这些群可以如下显式描述。按假设存在 $k$ 上分裂约化群 $G$，使得 $[N]^*\mathbb{G}|_{\GG_m}=G\times\GG_m$。覆盖 $[N]$ 的自同构群为 $\mu_N$。固定 $T\subset B\subset G$ 为分裂极大环面与 Borel 子群，以及 $G$ 的一个钉住 $\dagger$。

这些数据定义一个态射 $\sigma:\mu_N\to\operatorname{Aut}^\dagger(G)$，其中 $\operatorname{Aut}^\dagger(G)$ 是保持钉住的 $G$ 的自同构群。记 $[N]^*(G\times\GG_m)$ 为 $G\times\GG_m$ 沿覆盖 $[N]$ 的 Weil 限制。则 $\mu_N$ 通过 $\sigma$ 在 $G$ 上的作用与在覆盖 $[N]$ 上的作用而作用于这个 Weil 限制。我们仍把对角作用记为 $\sigma$。于是下降给出：
\[
\mathbb{G}|_{\PP^1\setminus\{0,\infty\}}\cong\big([N]^*(G\times\PP^1\setminus\{0,\infty\})\big)^\sigma
=(G\times\GG_m)/(\mu_N).
\]

\subsection{$0$ 与 $\infty$ 处的水平结构}

接下来需要描述 $\mathbb{G}|_{\PP^1\setminus\{0,\infty\}}$ 到 $\PP^1$ 的若干扩充。仍记 $[N]:\PP^1\to\PP^1$ 为由 $t\mapsto t^N$ 给出的映射。

群 $\mathbb{G}|_{\PP^1\setminus\{0,\infty\}}$ 将被扩充为由凸函数（在所有根上取 $0$）定义的（特殊）Bruhat--Tits 群 $\mathbb{G}$。由于我们假设 $\mathbb{G}|_{\PP^1\setminus\{0,\infty\}}$ 在驯扩张上分裂，这个群可以描述为群 $\big([N]^*(G\times\PP^1)\big)^\sigma$ 的连通分支。子群 $S\subset T|_{\PP^1\setminus\{0,\infty\}}\subset B|_{\PP^1\setminus\{0,\infty\}}\subset\mathbb{G}|_{\PP^1\setminus\{0,\infty\}}$ 定义闭子群 $S\subset\mathbb{T}\subset\mathbb{B}\subset\mathbb{G}$。

我们将需要在 $0$ 与 $\infty$ 处引入水平结构。它们对应于 $\mathbb{G}$ 的 Moy--Prasad 滤链的最初几步。关于这些群概形构造的自足论述可见 Yu 的预印本 [36]。

先考虑环面。对分裂环面 $T=\GG_m^r\times\PP^1$，记 $T(m_0,m_\infty)$ 为 $\PP^1$ 上光滑群概形，使得
\begin{gather*}
T(m_0,m_\infty)|_{\PP^1\setminus\{0,\infty\}}=\GG_m^r\times\PP^1\setminus\{0,\infty\},\\
T(m_0,m_\infty)(\mathcal O_{\PP^1,0})=\{g\in T(\mathcal O_{\PP^1,0})\mid g\equiv1\bmod t^{m_0}\},\\
T(m_0,m_\infty)(\mathcal O_{\PP^1,\infty})=\{g\in T(\mathcal O_{\PP^1,\infty})\mid g\equiv1\bmod s^{m_\infty}\}.
\end{gather*}
这也为诱导环面定义了滤链。对任意环面 $T$，取嵌入到诱导环面 $T\hookrightarrow I$，并通过拉回 $I$ 上的滤链来定义滤链。由 [36] 第 4 节，这个定义与所取嵌入无关；且由于我们假设 $T$ 在驯扩张上分裂，这些子群定义了连通群。

对约化群 $\mathbb{G}|_{\PP^1\setminus\{0,\infty\}}$，我们要为 $m_0,m_\infty\in\{0,1,2\}$ 定义模型 $\mathbb{G}(m_0,m_\infty)$，它们对应于 $0$ 处（相应地 $\infty$ 处）Iwahori 子群的 Moy--Prasad 滤链的第 $m_0$ 步（相应地第 $m_\infty$ 步）。

先假设 $\mathbb{G}=G\times\PP^1$ 是分裂半单群。像前面一样，固定 $S=T\subset B\subset G$ 为分裂极大环面、Borel 子群以及 $G$ 的钉住。记 $U\subset B$ 为幂幺子群。记 $\Phi=\Phi(G,S)$ 为根系，$\Phi^\pm\subset\Phi$ 为相对于所选 Borel 子群的正、负根之集。记 $\alpha_1,\dots,\alpha_r$ 为正单根，$\{\theta_j\}$ 为 $G$ 的最高根。最后，对每个根 $\alpha$，记 $U_\alpha$ 为 $G$ 中相应的根子群。

我们考虑 $G(k[[s]])$ 的如下有界子群：
\begin{itemize}[leftmargin=1.8em, itemsep=.2em]
\item $I(0):=\{g\in G(k[[s]])\mid g\bmod s\in B\}$ 是 Iwahori 子群。
\item $I(1):=\{g\in G(k[[s]])\mid g\bmod s\in U\}$ 是 $I(0)$ 的幂幺根。
\end{itemize}
为描述 $I(2)$，令 $f:\Phi\cup\{0\}\to\mathbb N$ 为由 $f(0)=1$ 以及
\[
f(\alpha)=
\begin{cases}
0 & \text{若 } \alpha\in\Phi^+\setminus\{\alpha_i\} \text{ 为正但非单根},\\
1 & \text{若 } \alpha=\alpha_i \text{ 是正单根},\\
1 & \text{若 } \alpha\in\Phi^-\setminus\{-\theta_j\},\\
2 & \text{若 } \alpha=-\theta_j \text{ 是某个最高根的负根}
\end{cases}
\]
定义的凹函数。$I(2)\subset G(k[[s]])$ 表示由凹函数 $f$ 定义的有界子群，即由 $\{u\in U_\alpha\mid u\equiv1\bmod s^{f(i)}\}$ 与 $\{g\in T(k[[s]])\mid g\equiv1\bmod s\}$ 生成的子群。

注意按定义 $I(1)/I(2)\cong\oplus_{\alpha\ \text{单仿射}}\GG_a$，同构于那些使 $\alpha$ 为单仿射根的根子群 $U_\alpha\subset G(k[[s]])$ 的直和。

类似地，我们用对偶 Borel 子群 $B^{\mathrm{opp}}$ 代替上述定义中的 Borel 子群，定义 $I(i)^{\mathrm{opp}}\subset G(k[[t]])$ 为相应的群。

\medskip
\noindent\textbf{例.} 对 $G=SL_n$，取 $S=T$ 为对角矩阵，$B$ 为上三角矩阵。则子群 $I(0)\subset SL_n(k[[s]])$ 由下三角项被 $s$ 整除的矩阵组成；$I(1)\subset I(0)$ 由对角项属于 $1+sk[[s]]$ 的矩阵组成。

$SL_n$ 关于单根的根子群由上对角项 $a_{i,i+1}$ 给出，而最长根之负根的 $U_{-\tilde\alpha}$ 由左下角的项给出。于是 $I(2)$ 由形如
\[
\begin{pmatrix}
1+sk[[s]] & sk[[s]] & k[[s]] & k[[s]]\\
sk[[s]] & 1+sk[[s]] & sk[[s]] & k[[s]]\\
sk[[s]] & sk[[s]] & 1+sk[[s]] & sk[[s]]\\
s^2k[[s]] & sk[[s]] & sk[[s]] & 1+k[[s]]
\end{pmatrix}
\]
的矩阵组成。把矩阵 $(a_{ij})$ 映到各项 $a_{i,i+1}$ 与 $a_{n,1}$ 的首项系数，得到同构 $I(1)/I(2)\cong\Aff^n$。

对一般分裂约化群 $\mathbb{G}=G\times\PP^1$，考虑导群 $G^{\mathrm{der}}$ 与中心的连通分支 $Z(G)^\circ$——后者是环面。暂时记 $I_{G^{\mathrm{der}}}(i)$、$I^{\mathrm{opp}}_{G^{\mathrm{der}}}(i)$ 为上面为半单群 $G^{\mathrm{der}}$ 定义的群，并定义
\[
I(i):=Z(G)^\circ(i)(k[[s]])\cdot I_{G^{\mathrm{der}}}(i)\subset G(k[[s]]),
\]
类似地 $I(i)^{\mathrm{opp}}=Z(G)^\circ(i,m_\infty)(k[[t]])\,I^{\mathrm{opp}}_{G^{\mathrm{der}}}(i)$。

对分裂约化群 $\mathbb{G}=G\times\PP^1$，群 $\mathbb{G}(m_0,m_\infty)$ 表示 Bruhat--Tits 群概形，使得
\begin{gather*}
\mathbb{G}(m_0,m_\infty)|_{\PP^1\setminus\{0,\infty\}}=G\times\GG_m,\\
\mathbb{G}(m_0,m_\infty)(\mathcal O_\infty)=I(m_\infty),\qquad
\mathbb{G}(m_0,m_\infty)(\mathcal O_0)=I(m_0)^{\mathrm{opp}}.
\end{gather*}
最后，对一般拟分裂群定义
\[
\mathbb{G}(m_0,m_\infty):=\Big(\big([N]^*\mathbb{G}(m_0,m_\infty)\big)^\sigma\Big)^\circ.
\]
注意对环面这并不给出新的定义。

我们将缩写：
\[
I(m_\infty):=\mathbb{G}(m_0,m_\infty)(\mathcal O_\infty),\qquad
I^-(m_0):=\mathbb{G}(m_0,m_\infty)(\PP^1-\{\infty\}).
\]
回忆 [32]：群 $I(m)$ 具有 $k$ 上（无限维）群概形的自然结构。

\subsection{仿射一般特征}

如前面所述，我们的构造依赖于 $I(1)$ 的一个特征的选择。称线性函数 $\varphi:I(1)/I(2)\to\Aff^1$ 为\emph{一般}的，若对任意单仿射根 $\alpha$，$\varphi$ 在 $U_\alpha$ 上的限制非平凡。全文固定这样一个一般 $\varphi$。

固定非平凡加性特征 $\psi:\FF_p\to\mathbb{Q}_\ell$，并把由 $\psi\circ\mathrm{Tr}_{k/\FF_p}$ 定义的 $k$ 上特征仍记为 $\psi$。用这个记号，特征 $\psi\circ\varphi$ 称为 $I(1)$ 的\emph{仿射一般特征}。

\subsection{主丛}

定义了群之后，我们需要收集关于 $G$-丛模叠 $\Bun_G$ 的几何的一些基本结果。所有这些对常群都是熟知的（例如可参看 [14] 的新近论述）。为把这些结果推广到我们的设定，我们依赖 [32] 与 [23]，其中解释了扭 loop 群的相应结果。需要指出的是，除 $\Bun_G$ 连通分支的计算外，所有结果都是关于 $\PP^1$ 上、在驯分歧覆盖 $[N]:\PP^1\to\PP^1$ 上分裂的群概形的特有问题。

首先回忆 [23] 中的一些结果与记号。记 $N$ 为 $T$ 的正规化子，$\widetilde W:=N(k((s)))/T(k[[s]])$ 为 Iwahori--Weyl 群。进一步记 $W_0:=N(k((s)))/T(k((s)))$ 为相对 Weyl 群。它同构于特殊纤维 $\mathbb{G}_{\infty,\mathrm{red}}$ 的约化商群的 Weyl 群（同处 Proposition 13），且
\[
\widetilde W\cong X_*(T)^{\pi_1(\PP^1\setminus\{0,\infty\})}\ltimes W_0.
\]
记 $W_a$ 为 $G(k((s)))$ 的根系（的）仿射 Weyl 群，它可以等同于 $G$ 的导群的单连通覆盖的 Iwahori--Weyl 群（[23] p.~196）。这是一个 Coxeter 群。于是 $\widetilde W\cong W_a\ltimes\Omega$，其中 $\Omega\cong X_*(Z(\ckG)^{\pi_1(\GG_m)})$ 是一个 alcove 的稳定子。

最后，对 $x\in\PP^1-\{\infty\}$，记 $\Gr_{G,x}$ 为仿射 Grassmannian（见 [32]），使得 $\Gr_{G,x}(k)=G(K_x)/G(\mathcal O_x)$。它也可以定义为参数化 $\PP^1$ 上 $G$-丛 $P$ 连同平凡化 $\phi:P|_{\PP^1\setminus\{x\}}\xrightarrow{\sim}\mathbb{G}_{\PP^1\setminus\{x\}}$ 的 ind-概形（例如 [24]）。

\begin{proposition}
假设基域 $k$ 或是有限域，或是代数闭域。设 $\mathbb{G}$ 为 $\PP^1$ 上 §1 中定义的拟分裂群概形。则下列结论成立：
\begin{enumerate}[label=(\arabic*), leftmargin=1.8em, itemsep=.2em]
\item 对任意 $x\in\PP^1$，典范映射 $\Gr_{G,x}\to\Bun_G$ 在 $\Bun_G$ 的光滑拓扑下局部有截面。此外，该映射在 $k$-点上本质满，即：$\PP^1$ 上任意 $G$-丛到 $\PP^1\setminus\{x\}$ 的限制都是平凡的。
\item $\pi_0(\Bun_G)\cong\pi_1\big([N]^*\mathbb{G}|_{\GG_m}\big)^{\pi_1(\GG_m)}\cong\Omega$。
\item $\PP^1$ 上每个 $G$-丛都可约化到 $T$。
\item （Birkhoff--Grothendieck 分解）
\[
\mathbb{G}(0,0)(k((s)))=\bigsqcup_{w\in\widetilde W}I^-(0)\cdot w\cdot I(0).
\]
\end{enumerate}
\end{proposition}

若 $\mathbb{G}=G\times\PP^1$ 是分裂群，本结果的证明可在 [14] 中找到。由于分裂群情形是我们的出发点，一般情形的证明将推迟到附录。对 $\gamma\in\Omega$，记 $\Bun^\gamma_{G(m,n)}$ 为 $\Bun_{G(m,n)}$ 的相应连通分支。

让我们收集该结果的一些推论。

\begin{corollary}
设 $n\in\{0,1,2\}$，$m\in\{0,1\}$。任意 $\gamma\in\Omega$ 定义同构 $\mathrm{Hk}_\gamma:\Bun_{G(m,n)}\to\Bun_{G(m,n)}$。它诱导连通分支之间的同构 $\Bun^0_{G(m,n)}\to\Bun^\gamma_{G(m,n)}$。特别地，$\Bun_{G(m,n)}$ 的所有连通分支互相同构。
\end{corollary}

\begin{proof}
由命题的第 (2) 部分，$\Bun_{G(m,n)}$ 的连通分支由 $\Omega\subset\widetilde W$ 编号。由于 $\Omega$ 是定义 $I(0)^{\mathrm{opp}}$ 的 alcove 的稳定子，任意 $\gamma\in\Omega$ 都正规化 $I(0)^{\mathrm{opp}}$ 及其幂幺根 $I(1)^{\mathrm{opp}}$。于是 $\Omega$ 通过右乘作用在 $\Gr_{G(m,n),0}=G(k((t)))/I(m)^{\mathrm{opp}}$ 上。这定义同构 $\mathrm{Hk}_\gamma:\Bun_{G(m,n)}\to\Bun_{G(m,n)}$。此外，[23] Lemma 14 蕴涵该映射在 $\Bun^0_{G(m,n)}$ 上的限制诱导同构 $\Bun^0_{G(m,n)}\to\Bun^\gamma_{G(m,n)}$。
\end{proof}

\begin{corollary}\label{cor:1.3}
\begin{enumerate}[label=(\arabic*), leftmargin=1.8em, itemsep=.2em]
\item 记 $\mathbb{G}(\PP^1)$ 为平凡 $G$-丛的自同构群。则包含 $\mathrm B\mathbb{G}(\PP^1)\to\Bun_G$ 是仿射开嵌入。
\item 记 $T(\PP^1)=H^0(\PP^1,T)=H^0(\PP^1,\mathbb{G}(0,0))$ 为平凡 $\mathbb{G}(0,0)$-丛的自同构群。平凡丛的包含定义仿射开嵌入 $\mathrm B(T(\PP^1))\hookrightarrow\Bun_{\mathbb{G}(0,0)}$。
\item 平凡 $\mathbb{G}(0,1)$-丛定义仿射开嵌入 $\operatorname{Spec}(k)\hookrightarrow\Bun^0_{\mathbb{G}(0,1)}$。
\item 把 $I(1)/I(2)$ 在 $\Bun_{\mathbb{G}(0,2)}$ 上的作用施于平凡 $\mathbb{G}(0,2)$-丛，得到典范映射 $j_0:I(1)/I(2)\hookrightarrow\Bun_{\mathbb{G}(0,2)}$。它是仿射开嵌入。
\item 对任意 $\gamma\in\Omega$，映射
\[
j_\gamma:=\mathrm{Hk}_\gamma\circ j_0:I(1)/I(2)\hookrightarrow\Bun_{\mathbb{G}(0,2)}
\]
是仿射开嵌入，称为\emph{大胞腔}。
\item 把 $T_0^{\mathrm{red}}\times I(1)/I(2)$ 在 $\Bun_{\mathbb{G}(1,2)}$ 上的作用施于平凡 $\mathbb{G}(1,2)$-丛，得到典范仿射嵌入：
\[
\tilde j_0:T_0^{\mathrm{red}}\times I(1)/I(2)\hookrightarrow\Bun_{\mathbb{G}(1,2)},\qquad
\tilde j_\gamma:=\mathrm{Hk}_\gamma\circ\tilde j_0:T_0^{\mathrm{red}}\times I(1)/I(2)\hookrightarrow\Bun_{\mathbb{G}(1,2)}.
\]
\end{enumerate}
\end{corollary}

\begin{proof}
先回忆为何本推论对 $GL_n$ 成立。对 (1)，注意 $\PP^1$ 上唯一上同调平凡的秩 $n$ 向量丛是 $\mathcal O(-1)^n$。因此 $\Bun_{GL_n}$ 上上同调行列式线丛之逆有一个零点恰在平凡丛上的截面。这证明了本情形的 (1)。其次回忆 $GL_n(0,0)$-丛就是向量丛连同 $0$ 与 $\infty$ 处的完整旗。记 $\mathrm{Mod}_{i,0}:\Bun^0_{GL_n(0,0)}\to\Bun^i_{GL_n(0,0)}$ 为沿 $0$ 处旗的第 $i$ 次上修改。其逆由第 $i$ 次下修改给出，我们记为 $\mathrm{Mod}_{-i,0}$。为证 (2)，记 $\mathrm{Mod}_{i,\infty}$ 为 $\infty$ 处的第 $i$ 次修改。平凡 $GL_n(0,0)$-丛是由 $\mathcal O^n$ 以及 $0$ 处纤维 $(\mathcal O^n)_0$ 的典范对偶旗 $(V_{i,0})_{i=0,\dots,n}$ 与 $\infty$ 处的 $(V_{i,\infty})_{i=0,\dots,n}$ 给出的丛。它是 $\PP^1$ 上唯一的、使复形 $E\to E_0/V_{i,0}\oplus E_\infty/V_{n-i,\infty}$ 对所有 $i$ 都上同调平凡的次数 $0$ 的 $GL_n(0,0)$-丛 $(E,V_{i,0},V_{i,\infty})$。因此平凡 $GL_n(0,0)$-丛在 $\Bun_{GL_n(0,0)}$ 中的包含又一次由线丛截面的非消没定义。

对一般 $G$，取忠实表示 $\rho:\mathbb{G}\to GL_n\times\PP^1$。它定义映射 $\Bun_G\to\Bun_{GL_n}$。Birkhoff--Grothendieck 分解蕴涵：$G$-丛平凡当且仅当相伴的 $GL_n$-丛平凡。此外，为验证 $0$、$\infty$ 处的约化到 $B$ 是互相相反的，也只需在诱导的 $GL_n$-丛上验证。这证明了 (1) 与 (2)。

(3) 由 (2) 以及命题 1.1 证明中的同构 $H^0(\PP^1,T)\cong T_0^{\mathrm{red}}$ 推出。(4) 由 (3) 推出，因为映射 $\Bun_{\mathbb{G}(0,2)}\to\Bun_{\mathbb{G}(0,1)}$ 是 $I(1)/I(2)$-torsor。由推论 1.2，(5) 由 (4) 推出。最后，(6) 由 (5) 推出，因为 $\Bun_{\mathbb{G}(1,2)}\to\Bun_{\mathbb{G}(0,2)}$ 是 $T_0^{\mathrm{red}}$-torsor。
\end{proof}
